\documentclass[11pt,a4paper]{article}

\usepackage[utf8]{inputenc}
\usepackage[T1]{fontenc}
\usepackage[french]{babel}
\usepackage[margin=2.4cm]{geometry}
\usepackage{booktabs}
\usepackage{enumitem}
\usepackage{xcolor}
\usepackage{hyperref}
\hypersetup{colorlinks=true, linkcolor=blue!50!black, urlcolor=blue!50!black}

\newcommand{\code}[1]{\texttt{#1}}

% Boîtes légères (xcolor seulement, portables)
\newenvironment{say}{\par\medskip\noindent\textcolor{blue!45!black}{\rule[-1pt]{3pt}{12pt}}\
  \textbf{\textcolor{blue!45!black}{À dire}}\par\nopagebreak\smallskip\noindent\ignorespaces}{\par\medskip}
\newenvironment{why}{\par\medskip\noindent\textcolor{orange!60!black}{\rule[-1pt]{3pt}{12pt}}\
  \textbf{\textcolor{orange!60!black}{Pourquoi ce choix}}\par\nopagebreak\smallskip\noindent\ignorespaces}{\par\medskip}
\newenvironment{qa}{\par\medskip\noindent\textcolor{gray!55!black}{\rule[-1pt]{3pt}{12pt}}\
  \textbf{\textcolor{gray!55!black}{Questions probables}}\par\nopagebreak\smallskip\noindent\ignorespaces}{\par\medskip}

\title{\textbf{Absentéisme au travail : séries temporelles \& classification}\\[2mm]
\large Notes pour l'oral --- Time Series Analysis}
\author{Groupe : \underline{\hspace{4cm}} \quad / \quad \underline{\hspace{4cm}}}
\date{Soutenance --- juin 2026}

\begin{document}
\maketitle

\begin{abstract}
\noindent
Document de préparation à l'oral. Il résume le notebook \code{main.ipynb}, justifie chaque
choix de modèle / fonction, indique \textbf{qui dit quoi} (2 présentateurs) et anticipe les
questions du jury. Couleurs : \colorbox{blue!8}{bleu = à dire}, \colorbox{orange!10}{orange =
justification}, \colorbox{gray!10}{gris = Q/R}.
\end{abstract}

%==================================================================
\section{Contexte et problématique}

\begin{say}
Jeu de données \textbf{Absenteeism at Work} (UCI \#445) : 740 enregistrements d'absences chez
un transporteur brésilien (2007--2010), 21 variables (motif médical, jour, charge de travail,
distance, profil de l'employé\dots) et une cible : \emph{heures d'absence}.

Deux questions :
\begin{enumerate}[nosep]
  \item \textbf{Prévision} : peut-on anticiper le volume mensuel d'heures d'absence (planification RH) ?
  \item \textbf{Explication} : peut-on prédire qu'une absence sera \emph{longue} (> 8 h) à partir de ses caractéristiques ?
\end{enumerate}
\end{say}

%==================================================================
\section{Répartition orale (2 présentateurs)}

\begin{center}
\begin{tabular}{@{}p{1.3cm}p{8.4cm}p{2.6cm}@{}}
\toprule
& \textbf{Contenu} & \textbf{Durée indic.} \\
\midrule
\textbf{P1} & Intro + problématique, \S2 \emph{Analyse qualitative}, \S3 \emph{Forecasting} (partie séries temporelles) & $\approx$ 6--7 min \\
\textbf{P2} & \S4 \emph{Effet du jour}, \S5 \emph{Classification}, conclusion \& perspectives (partie explicative) & $\approx$ 6--7 min \\
\bottomrule
\end{tabular}
\end{center}

\noindent\emph{Découpage naturel « séries temporelles / apprentissage explicatif ». Chacun doit
aussi pouvoir répondre sur la partie de l'autre (cf. grille).}

%==================================================================
\section{Couverture de la grille de compétences}

\begin{center}
\begin{tabular}{@{}p{5.2cm}p{6.4cm}p{2.4cm}@{}}
\toprule
\textbf{Compétence (AAV)} & \textbf{Où / comment} & \textbf{Présent.} \\
\midrule
Étude statistique qualitative (description, graphiques) & \code{describe}, séries, moyenne mobile, sous-séries annuelles & P1 \\
Étude heuristique (stationnarité, saisonnalité, tendance) & Décomposition additive/multiplicative + STL, ADF+KPSS, différenciation, ACF/PACF & P1 \\
Forecasting : choix \& justification & Baselines, AR(p), MA(q), ARIMA(p,d,q), sélection AIC/BIC & P1 \\
Forecasting : mise en œuvre & Split chronologique, estimation OLS/MLE, préprocessing \code{d=0} justifié & P1 \\
Forecasting : évaluation/critique & MAE/RMSE vs baselines, Ljung-Box (résidus), limites & P1 \\
Forecasting : présentation & Prévision test + IC 95\,\%, prévision 12 mois & P1 \\
Classification : choix \& justification & Dummy / LogReg / Random Forest, métriques adaptées au déséquilibre & P2 \\
Classification : mise en œuvre & Pipeline \code{ColumnTransformer}, CV 5-fold stratifiée, \code{class\_weight} & P2 \\
Classification : évaluation/critique & ROC/PR-AUC/F1, matrice de confusion, importances, limites & P2 \\
Problématique \& recontextualisation & Intro + conclusion (RH) & P1 / P2 \\
\bottomrule
\end{tabular}
\end{center}

%==================================================================
\section*{\textcolor{blue!50!black}{PRÉSENTATEUR 1 --- Séries temporelles}}
\addcontentsline{toc}{section}{Présentateur 1}

\subsection{Construction de la série (\S2)}
\begin{say}
Le dataset n'a \textbf{pas de colonne date}, mais les lignes sont ordonnées chronologiquement.
On reconstruit un index mensuel (on incrémente l'année quand le mois « repart en arrière »),
puis on \textbf{agrège le total d'heures par mois} : série mensuelle de \textbf{37 points}
(juil.~2007 $\to$ juil.~2010).
\end{say}
\begin{why}
Somme (et non moyenne) par mois car la quantité d'intérêt RH est le \emph{volume} total d'heures
perdues. \code{asfreq('MS')} fixe une fréquence mensuelle régulière, indispensable aux outils TSA.
\end{why}

\subsection{Description et tendance}
\begin{say}
\code{describe} (min/max/moyenne), tracé de la série brute, puis \textbf{moyenne mobile} centrée
(fenêtres 3 et 12 mois) et \textbf{sous-séries annuelles} (une courbe par année).
\end{say}
\begin{why}
Moyenne mobile (Ch.3) : moyenner sur \emph{un cycle complet} (12 mois) annule la composante
saisonnière et fait ressortir la tendance ; la fenêtre 3 mois garde le détail trimestriel.
\end{why}

\subsection{Décomposition saisonnière}
\begin{say}
Décomposition \textbf{additive} ($X=T+S+R$) et \textbf{multiplicative} ($X=T\cdot S\cdot R$) via
\code{seasonal\_decompose}, puis \textbf{STL} (Loess).
\end{say}
\begin{why}
Additif si l'amplitude saisonnière est stable ; multiplicatif si elle croît avec le niveau.
\textbf{STL} (Ch.3) lisse par régressions locales (Loess) et l'option \code{robust=True} réduit
l'influence des pics (ex.\ COVID) : trend plus propre que la moyenne mobile.
\end{why}

\subsection{Stationnarité}
\begin{say}
Moyenne / écart-type mobiles (12 mois) + deux tests : \textbf{ADF} ($H_0$ : racine unitaire) et
\textbf{KPSS} ($H_0$ : stationnaire). Résultats : \textbf{ADF $p=0{,}000$} et \textbf{KPSS
$p=0{,}100$} $\Rightarrow$ la série est \textbf{stationnaire}.
\end{say}
\begin{why}
On utilise les \emph{deux} tests ensemble (Ch.3) car ils ont des hypothèses \textbf{opposées} :
l'accord (ADF rejette \emph{et} KPSS ne rejette pas) est bien plus convaincant qu'un seul test.
Conséquence directe : \textbf{ordre d'intégration $d=0$}.
\end{why}

\subsection{Différenciation et ACF/PACF}
\begin{say}
On montre quand même la différenciation $\nabla = (1-B)$ et on retrace ADF/KPSS, puis l'ACF et la
PACF (brute vs différenciée).
\end{say}
\begin{why}
$\nabla$ (Ch.3) retirerait une tendance, mais la série étant déjà stationnaire, différencier
serait de la \textbf{sur-différenciation}. L'ACF décroît vite et il n'y a \textbf{pas de pic au
lag 12} $\Rightarrow$ pas de saisonnalité annuelle marquée (justifie l'absence de SARIMA).
\end{why}

\subsection{Forecasting (\S3)}
\begin{say}
Cadre : un bon modèle laisse un résidu = \textbf{bruit blanc}. Découpage \textbf{chronologique}
train/test = \textbf{28 / 9 mois}. Métriques \textbf{MAE} et \textbf{RMSE}. On compare, du plus
simple au plus riche :
\begin{itemize}[nosep]
  \item \textbf{Baselines} : moyenne (RMSE 98{,}5), naïf (114{,}1), naïf saisonnier (108{,}3).
  \item \textbf{AR(2) par OLS} (recette Ch.2) : \textbf{RMSE 95{,}9} --- le meilleur.
  \item \textbf{MA(2)} et \textbf{ARIMA(0,0,2)} par MLE : RMSE 97{,}7.
\end{itemize}
\end{say}
\begin{why}
\textbf{Split chronologique} (et non aléatoire) : on prédit le futur à partir du passé, sinon
fuite temporelle. \textbf{AR(p)} suit littéralement les \emph{Steps 1--3} du Ch.2 : (1) $p$ = dernier
lag significatif de la \textbf{PACF} $\Rightarrow p=2$ ; (2) matrice de retards ; (3) régression
\textbf{OLS} pour $\Phi(B)$ ($\Phi\approx[-0{,}54;\ 0{,}05]$) --- « un AR pur s'estime par OLS ».
\textbf{MA(q)} : $q$ lu sur la \emph{coupure} de l'ACF. \textbf{ARIMA(p,d,q)} : grille balayée,
ordre retenu par \textbf{BIC} (plus parcimonieux et cohérent que l'AIC) ; ici $d=0$ donc ARIMA
$=$ ARMA. MA/ARMA s'estiment par \textbf{MLE (Kalman)}, contrairement à l'AR.
\end{why}
\begin{say}
\textbf{Diagnostic} : ACF des résidus plate + \textbf{Ljung-Box} $p=0{,}94$ (lag 6) et $0{,}87$
(lag 12) $\Rightarrow$ résidus compatibles avec un \textbf{bruit blanc}. On présente la prévision
sur le test avec \textbf{IC 95\,\%} puis une prévision \textbf{12 mois} sur série complète.
\end{say}
\begin{why}
\textbf{Conclusion \S3 à assumer} : les modèles ne battent que \emph{marginalement} la baseline
(95{,}9 vs 98{,}5) et les résidus sont \emph{déjà} du bruit blanc $\Rightarrow$ la série mensuelle
(37 points) porte \textbf{peu de dynamique exploitable}. C'est un \emph{résultat}, pas un échec :
il justifie le pivot vers l'approche explicative (P2).
\end{why}

%==================================================================
\section*{\textcolor{blue!50!black}{PRÉSENTATEUR 2 --- Analyse explicative}}
\addcontentsline{toc}{section}{Présentateur 2}

\subsection{Effet du jour de la semaine (\S4)}
\begin{say}
On change d'angle : le \code{Day of the week} est la seule dimension temporelle fine fiable.
Statistiques par jour : \textbf{lundi} le plus lourd (9{,}2 h en moyenne), \textbf{jeudi} le plus
léger (4{,}4 h). Tests : \textbf{Kruskal-Wallis} $H=10{,}7,\ p=0{,}030$ (les heures dépendent du
jour) ; \textbf{$\chi^2$} d'indépendance jour $\times$ « absence longue » : $\chi^2=11{,}52$,
ddl 4, $p=0{,}021$ (dépendance significative).
\end{say}
\begin{why}
\textbf{Kruskal-Wallis} (non paramétrique) plutôt qu'une ANOVA : la durée d'absence est très
\textbf{asymétrique} (non gaussienne), Kruskal ne suppose pas la normalité. \textbf{$\chi^2$} :
adapté à deux variables \emph{catégorielles} (jour $\times$ binaire). But : montrer que le jour est
un \textbf{prédicteur pertinent} avant la classification.
\end{why}

\subsection{Classification : prédire une absence longue (\S5)}
\begin{say}
Cible binaire \code{absence\_longue} = (durée > 8 h). Classe \textbf{rare} : \textbf{63/737 =
8{,}5\,\%}. On évalue par \textbf{ROC-AUC}, \textbf{PR-AUC} et \textbf{F1} --- \textbf{pas
l'accuracy}.
\end{say}
\begin{why}
Avec 8{,}5\,\% de positifs, un modèle qui dit toujours « non » a déjà 91{,}5\,\% d'accuracy : la
métrique serait trompeuse. \textbf{PR-AUC} et \textbf{F1} sont pertinentes pour les classes rares ;
\textbf{ROC-AUC} mesure le pouvoir de séparation indépendamment du seuil.
\end{why}
\begin{say}
\textbf{Préprocessing} : \code{ColumnTransformer} --- one-hot pour les catégorielles (motif, jour,
saison, éducation), standardisation pour les numériques (charge, distance, âge\dots), binaires
laissées telles quelles. Le tout dans un \code{Pipeline}. \textbf{Modèles} : Dummy (baseline),
\textbf{Régression logistique}, \textbf{Random Forest}, avec \code{class\_weight='balanced'}.
Évaluation par \textbf{CV 5-fold stratifiée}.
\end{say}
\begin{why}
\textbf{Pipeline} : préprocessing appris \emph{dans} la CV $\Rightarrow$ pas de \emph{fuite de
données}. \textbf{Standardisation} utile à la régression logistique (échelles comparables) ;
\textbf{one-hot} car les codes (motif, jour) sont catégoriels, pas ordinaux.
\code{class\_weight='balanced'} compense le déséquilibre. \textbf{CV stratifiée} : conserve les
8{,}5\,\% de positifs dans chaque pli, plus robuste qu'un seul split vu le faible nombre de positifs.
\end{why}
\begin{say}
\textbf{Résultats} (CV ROC-AUC) : Random Forest \textbf{0{,}84}, Régression logistique
\textbf{0{,}81}, baseline 0{,}50. Sur le \emph{test}, scores plus bas (RF 0{,}75) car seulement
$\sim$16 positifs $\Rightarrow$ forte variance : \textbf{la CV est la mesure fiable}. Variables
les plus discriminantes : \textbf{motif} d'absence et \textbf{charge de travail} ; le jour de la
semaine pèse $\approx$ \textbf{10\,\%} (confirme le \S4, en facteur secondaire).
\end{say}

\subsection{Conclusion générale \& perspectives}
\begin{say}
\textbf{Recontextualisation RH} : (1) la prévision mensuelle agrégée est peu fiable (série $\approx$
bruit) $\Rightarrow$ ne pas surinvestir dans un modèle ARIMA pour le pilotage ; (2) en revanche on
\emph{explique} les absences longues (motif, charge, lundi) $\Rightarrow$ leviers concrets
(lissage de charge, vigilance lundi). \textbf{Limites} : 37 points pour la série, classe rare en
classification. \textbf{Pistes} : rééchantillonnage (SMOTE) / ajustement du seuil, gradient
boosting, exploitation de l'historique par employé (\code{ID}).
\end{say}

%==================================================================
\section{Questions probables du jury (Q/R)}

\begin{qa}
\textbf{Pourquoi $d=0$ dans l'ARIMA ?} La série est stationnaire (ADF $p=0{,}000$ \emph{et} KPSS
$p=0{,}100$) ; différencier serait de la sur-différenciation. ARIMA(p,0,q) $=$ ARMA(p,q).

\textbf{Pourquoi pas de SARIMA ?} Pas de pic d'ACF aux lags 12, 24 $\Rightarrow$ pas de
saisonnalité annuelle nette à modéliser ; on l'écarte de façon justifiée.

\textbf{AR par OLS, est-ce correct ?} Oui : pour un AR \emph{pur}, l'estimation OLS sur la matrice
de retards est valide (le cours le dit). MA/ARMA, eux, exigent une MLE (Kalman). Nuance : notre OLS
inclut une constante (la série n'est pas centrée en 0), ce qui n'apparaît pas dans la formule
« nue » de la slide mais capte la moyenne.

\textbf{ACF vs PACF ?} ACF = corrélation avec le retard $h$ (sert à lire $q$ d'un MA, qui coupe à
$q$). PACF = corrélation \emph{directe} au retard $h$, effets intermédiaires retirés (sert à lire
$p$ d'un AR, qui coupe à $p$).

\textbf{AIC ou BIC ?} On choisit par BIC (pénalité $(p+q)\ln T$, plus forte) : plus parcimonieux et
\emph{consistant} (retrouve le vrai ordre asymptotiquement) ; l'AIC sur-ajuste en petit échantillon.

\textbf{Qu'est-ce qu'un bruit blanc et pourquoi Ljung-Box ?} Bruit blanc : moyenne nulle, variance
constante, non corrélé ($\rho(h)=\mathbf{1}_{h=0}$). Si les résidus en sont un, le modèle a capté
toute la structure. Ljung-Box teste l'absence d'autocorrélation jusqu'au lag $k$ ($p>0{,}05$ = OK).

\textbf{Split chronologique, pourquoi ?} Un split aléatoire utiliserait le futur pour prédire le
passé (fuite) ; en série temporelle on entraîne sur le passé, on teste sur le futur.

\textbf{Additif ou multiplicatif ?} Additif si l'amplitude saisonnière est stable, multiplicatif si
elle croît avec le niveau (souvent réglé en passant au log). On montre les deux.

\textbf{Pourquoi pas l'accuracy en classification ?} Classe à 8{,}5\,\% : « toujours négatif » fait
91{,}5\,\% d'accuracy sans rien prédire. D'où ROC-AUC / PR-AUC / F1.

\textbf{Pourquoi la CV plutôt que le test ?} $\sim$16 positifs en test $\Rightarrow$ score très
variable ; la CV 5-fold stratifiée moyenne sur plusieurs découpes, plus fiable.

\textbf{Risque de fuite de données ?} Non : préprocessing (one-hot, scaler) encapsulé dans le
\code{Pipeline}, donc ré-appris sur chaque pli d'entraînement uniquement.

\textbf{Le motif ne prédit-il pas « mécaniquement » la durée ?} En partie : certains motifs sont
quasi-déterministes. La vraie valeur ajoutée du modèle est sur les cas ambigus ; à mentionner comme
limite.
\end{qa}

%==================================================================
\section{Conseils de présentation}
\begin{itemize}[nosep]
  \item \textbf{Temps} : viser 12--14 min à deux, garder 1 figure forte par section (série + décomposition ; recap des scores ; ROC + importances).
  \item \textbf{Transitions} : P1 termine sur « la série est $\approx$ un bruit blanc » $\Rightarrow$ P2 enchaîne « donc on explique plutôt que prédire ».
  \item \textbf{Chiffres à retenir} : 37 points ; ADF 0{,}000 / KPSS 0{,}100 ; AR(2) RMSE 95{,}9 ; 8{,}5\,\% de positifs ; CV ROC-AUC 0{,}84.
  \item \textbf{Niveau de langage} : dire « stationnaire », « autocorrélation », « parcimonie », « déséquilibre de classes » ; éviter le familier.
\end{itemize}

\end{document}
